\section{Problem Description}

To address the research questions, we will first conduct a thorough exploratory data analysis to extract meaningful insights from the dataset. Based on these initial findings, we will formulate hypotheses and develop preliminary claims regarding the differences between Disney+ and Netflix. These claims will then be subjected to rigorous statistical validation using appropriate methods, such as the t-test, to ensure their reliability and accuracy. Through this process, we aim to draw robust conclusions and provide answers to the proposed research questions.

\subsection{Dataset Overview}

This section provides a foundational understanding of the dataset, outlining its structure and key variables to set the stage for further analysis.

\subsubsection{Data Source and Collection}

The dataset comprises data that was scraped from multiple streaming platforms, including Netflix, Disney+, Hulu, and Prime Video. The scraping process provided a comprehensive list of movies available on these platforms, capturing key metadata. This method ensured a systematic and extensive collection of data from publicly accessible information.

\subsubsection{Dataset Variables and Their Characteristics}

The dataset contains \textbf{9515 records} and represents movies from major streaming platforms. This dataset consists of 11 columns with various data types. Most columns, such as \texttt{ID}, \texttt{Year}, and platform indicators (\texttt{Netflix}, \texttt{Hulu}, etc.), are integers (\texttt{int64}). The \texttt{Title} column stores movie names as objects and generally does not require processing. However, the \texttt{Age} and \texttt{Rotten Tomatoes} columns, also stored as objects, need to be converted to appropriate formats for analysis and statistical testing. \textbf{Figure~\ref{fig:datatable}} in the appendix shows the first few rows of the dataset.

\subsubsection{Missing Values}

The dataset contains missing values in two key columns:

\begin{itemize}
    \item \textbf{Age}: 
    \begin{itemize}
        \item \textbf{Missing Values}: 4,177 entries
        \item \textbf{Percentage}: 43.9\%
    \end{itemize}
    This is a significant proportion of the dataset, suggesting that nearly half of the records lack age ratings. Since this variable is crucial for analyzing age suitability, appropriate handling of these missing values is necessary. Options include imputation (e.g., replacing missing values with the most frequent category) or excluding these rows, depending on the analysis requirements.

    \item \textbf{Rotten Tomatoes}:
    \begin{itemize}
        \item \textbf{Missing Values}: 7 entries
        \item \textbf{Percentage}: 0.07\%
    \end{itemize}
    The number of missing values here is minimal and unlikely to affect the overall analysis significantly. These can either be imputed using the mean or median score or removed without substantial impact.
\end{itemize}
